\chapter{绪论}

\section{研究背景及意义}

气候变化已经成为全球能源系统转型必须面对的长期约束。为控制温室气体排放，《巴黎协定》提出将全球平均气温升幅控制在较工业化前水平明显低于 \(2\,^\circ\mathrm{C}\) 的范围内，并努力限制在 \(1.5\,^\circ\mathrm{C}\) 以内\cite{unfccc_paris_2015}。中国作为全球重要能源消费国和制造业大国，提出力争于 2030 年前实现二氧化碳排放达峰、2060 年前实现碳中和\cite{china_ndc_2021}。在这一目标牵引下，电力系统需要在保障供电安全的同时持续降低单位电量碳排放，能源结构也正由以煤电为主体的传统形态向高比例可再生能源与多元调节资源协同的新型电力系统转变。

风电、光伏等可再生能源具有清洁低碳优势，但其出力受气象条件影响显著，存在波动性、间歇性和预测误差。随着新能源装机占比提升，电力系统对灵活调节电源、快速爬坡能力和备用容量的需求同步增加。燃气轮机及燃气-蒸汽联合循环机组具有启停速度快、负荷调节范围宽、单位发电碳排放低于常规燃煤机组、污染物排放控制相对容易等特点，能够在新能源消纳、电网调峰、冷热电联供和应急备用等场景中发挥重要作用。国际能源署等机构也指出，在高比例可再生能源系统中，灵活可调的低排放电源和储能、需求响应等资源共同构成系统安全运行的重要支撑\cite{iea_weo_2023}。因此，在较长时期内，燃气机组并不是简单的过渡性设备，而是连接化石能源清洁利用、可再生能源大规模接入和电力系统可靠运行的关键装备。

燃气-蒸汽联合循环机组通过燃气轮机布雷顿循环与蒸汽朗肯循环的耦合，实现燃料化学能的梯级利用，具有较高的热效率和较好的工程成熟度。典型联合循环系统中，燃气轮机排出的高温烟气进入余热锅炉产生蒸汽，并进一步推动汽轮机做功，从而提高系统总出力和燃料利用效率。其基本流程如图 \ref{fig:chapter1-ccpp-schematic} 所示。与单一循环相比，联合循环系统内部包含压气机、燃烧室、燃气透平、余热锅炉、汽轮机、凝汽器、给水系统和热网抽汽等多个相互耦合的设备，热力过程跨越气体侧和汽水侧，运行性能受环境条件、燃料组成、负荷水平、设备退化和控制策略等多重因素共同影响。

\begin{figure}[!htb]
  \centering
  \includegraphics[width=0.86\textwidth]{chapter1/燃气蒸汽联合循环示意图.png}
  \caption{燃气-蒸汽联合循环系统示意图}
  \label{fig:chapter1-ccpp-schematic}
\end{figure}

在工程运行中，燃气机组往往长期处于偏离设计工况的状态。一方面，电网调峰和供热需求使机组频繁经历部分负荷、启停和变负荷运行；另一方面，环境温度、压力、湿度和燃料组分的变化会改变压气机进气条件、燃烧过程和余热回收特性；此外，设备长期运行还会引起压气机污垢、透平通流能力变化、余热锅炉换热性能下降、汽轮机通流部分结垢等退化现象。这些因素会使压比、空气流量、透平排气温度、蒸汽参数、真空度和机组效率等关键指标偏离设计值，导致运行状态、能量分配和经济性发生变化。若仍以设计工况或静态经验指标评价机组性能，容易掩盖工况变化与设备退化的差异，也难以为高效可靠运行提供可靠依据。

因此，有必要面向真实运行数据重新评估燃气-蒸汽联合循环机组的运行状态，建立能够反映多工质、多设备和变工况特征的热力分析方法。通过对实测数据进行工质状态构造、设备能量分析和关键参数影响分析，可以识别机组在实际运行中的能量转换规律和主要损失环节，区分环境扰动、负荷变化和设备性能变化对效率的影响，为运行人员开展参数调整、负荷分配、维护决策和节能降耗提供定量依据。这不仅有助于提高单台机组的运行经济性，也能够在更大尺度上支撑新型电力系统低碳、安全和灵活运行。

\section{国内外研究现状}

燃气-蒸汽联合循环机组性能研究首先依赖于对系统热力过程的合理描述。现有研究通常以压气机、燃烧室、燃气透平、余热锅炉、汽轮机和凝汽器等关键设备为对象，基于质量守恒、能量守恒、部件特性关系和换热关系建立系统级热力模型。早期气路分析方法以燃气轮机部件性能偏差识别为核心，为后续燃气轮机状态评估和性能诊断奠定了基础\cite{urban1973gpa,stamatis2011gpa,fentaye2019review}；压气机和燃气透平特性建模则通过通用特性、变工况修正和参数校正等方法，为部件级性能计算提供支撑\cite{zhang1996compressor,lu1996turbine,volponi1999corrections}。面向工程应用的研究进一步将机理模型、现场测量和健康监测结合起来，用于燃气轮机性能评估、诊断和预测性维护\cite{tahan2017health,hanachi2018survey}。对于联合循环机组而言，这类模型能够把燃气轮机侧的布雷顿循环与汽水侧的朗肯循环联系起来，适合分析负荷变化、环境条件变化和设备边界条件变化对机组效率、出力和热耗率的影响。近年来，面向实际机组的模块化热力建模和运行数据修正逐渐成为提高模型工程适用性的重要方法\cite{wang2023hrsg,tsoutsanis2014map,ying2023diagnostic}。

在性能影响因素方面，国内外研究均指出联合循环机组效率不是单一参数决定的结果，而是负荷、环境条件、燃气轮机运行参数和蒸汽循环边界条件共同作用的结果。对于燃气轮机侧，环境温度、环境压力、进气流量、燃料流量、压气机压比和透平入口温度会影响压气机耗功、燃烧室出口状态、透平输出功和排气余热；对于蒸汽循环侧，余热锅炉换热过程、主蒸汽流量与压力、凝汽器真空度以及供热抽汽条件会影响汽轮机出力和全厂能量利用水平。已有研究从部分负荷特性、不同控制策略和环境温度影响等角度分析了联合循环变工况性能，指出燃气轮机与蒸汽循环之间存在明显耦合，系统性能需要从全厂角度评价\cite{kim2004partload,hyun2018dual,yang2019inlet,zhang2018chp,li2018backpressure}。同时，多压力联合循环、余热锅炉参数和蒸汽压力等因素也会影响烟气余热回收效率和底循环出力\cite{ibrahim2014multipressure,yang2018hrsg,liying2006steam}。因此，机组在部分负荷和变环境条件下常常偏离设计工况，需要从燃气轮机、余热锅炉和汽轮机之间的能量传递关系出发分析参数影响，而不能只考察单一设备或单一效率指标。本文后续对环境参数、燃机出力、汽机出力、给水压力、主蒸汽流量和供热量等变量的分析，正是沿着这一研究方向展开。

随着燃气机组承担更多调峰和灵活运行任务，长期宽工况运行下的状态评估逐渐成为研究重点。燃气轮机在长期运行中会受到积垢、腐蚀、磨损和热疲劳等因素影响，导致部件效率、通流能力和排气状态逐渐偏离设计基准\cite{diakunchak1992deterioration,kurz2001degradation}。联合循环系统中，燃气轮机、余热锅炉和汽轮机性能变化还会通过排气温度、烟气流量和蒸汽参数相互传递，从而影响全厂效率和负荷分配；余热锅炉换热面结垢、汽轮机通流部分退化等问题也会改变底循环吸热和做功过程\cite{marner1984fouling,zwebek2003degradation,massardo2009hrsg}。已有退化和性能诊断研究表明，若不能区分环境扰动、工况变化与设备性能变化，容易将正常变工况影响误判为设备退化\cite{li2009prognostic,tahan2017health}，因此需要通过实测数据重新构造机组运行状态，进行设备性能评估。

实际生产数据还存在测点不完整、关键参数难以直接测量和设备状态难以准确标注等问题。现场性能试验和能量监测研究指出，流量、温度、压力等测量量的不确定性会直接影响燃气轮机和联合循环性能计算结果\cite{pinelli2000uncertainties,carstens2018uncertainty}；同时，燃烧室出口温度、部分流量和湿蒸汽状态等关键参数往往难以直接测得，设备特性曲线和内部结构参数也可能因商业保密或现场条件限制而不完备\cite{colera2019thermodynamic,vonmoll2014egt,vemulapalli2022softsensor,tian2018softmeasurement}。这些问题说明，在真实联合循环机组上开展性能分析时，应尽量利用已有测点建立统一的工质状态和设备平衡关系，并通过数据一致性检查提高分析结果的可信度。

除能量分析外，㶲分析也被广泛用于联合循环系统性能评价。㶲分析以热力学第二定律为基础，能够进一步揭示不同设备中的不可逆损失和能量品质下降位置。已有研究将㶲分析与参数优化、多目标优化和环境经济评价相结合，用于识别燃烧室、余热锅炉、透平等设备的主要损失来源，并分析压比、透平入口温度、蒸汽参数等变量对系统性能的影响\cite{kaviri2012exergy,ahmadi2011exergo}。这类研究说明，仅从能量守恒角度分析联合循环机组仍不够充分，还需要结合㶲效率和㶲损失评价不同能量流的可用能水平。本文在实测数据基础上建立了包含㶲计算的统一流股模型，为后续进一步开展㶲分析预留了接口。受数据条件和研究周期限制，本文以能量分析为主线，重点从热力学第一定律角度刻画真实机组的能量利用过程。

\section{研究内容与章节安排}

本文以某燃气-蒸汽联合循环机组为研究对象，所用数据来自该机组 2025 年 4 月至 12 月的实际运行记录。数据覆盖燃气轮机、余热锅炉、汽轮机、凝汽器和热网抽汽等主要系统，包含环境参数、工质参数、部件参数等相关变量。由于该数据来源于生产现场，其工况分布并非理想设计点附近的少量稳态样本，而是包含季节变化、负荷调节、供热需求变化和设备状态波动的长期运行数据。因此，本文研究具有较强的工程真实性，也要求模型能够处理变工况、多变量和多设备耦合问题。

围绕上述对象，本文主要开展以下研究工作。第一，建立联合循环机组多工质流股建模方法。针对燃料气、湿空气、烟气、水和水蒸气等不同工质，构造统一的流股对象，计算温度、压力、质量流量、比焓、比熵、能量流和㶲流等状态量，为后续设备模型提供一致的数据接口。第二，建立燃气轮机、余热锅炉和汽轮机等关键设备的热力分析模型。基于热力学定律计算各设备的能量输入输出、功率分配、换热特性、能量损失和效率指标，从设备级揭示实际运行中的性能变化规律。第三，结合真实运行数据分析联合循环机组关键运行参数对性能的影响。通过对环境温度、环境压力、环境湿度、燃机出力、汽机出力、供热量、给水压力、主蒸汽流量和透平排气温度等变量进行相关性和散点规律分析，识别影响整机效率的主要因素，为运行诊断和性能评估提供依据。本文的研究建立了基于燃气-蒸汽联合循环机组的热力分析模型，为机组运行状态评估和效率分析奠定基础。

全文章节安排如下。

第 1 章为绪论，介绍低碳能源转型背景下燃气-蒸汽联合循环机组的重要作用，总结国内外关于性能评估和退化识别的研究现状，并说明本文研究内容与数据来源。

第 2 章介绍联合循环机组工质建模方法，建立燃料气、湿空气、烟气、水和水蒸气的物性计算与统一流股描述方法，并给出由实测数据构造全厂流股对象的过程。

第 3 章建立燃气轮机侧性能分析模型，围绕压气机、燃烧室和燃气透平等部件开展能量分析，研究燃气轮机在实际运行数据下的状态变化。

第 4 章建立余热锅炉和汽轮机侧热力分析模型，分析烟气余热利用、汽水侧吸热、汽轮机出力分配和相关效率指标。

第 5 章基于实际运行数据开展机组性能影响因素分析，结合相关性矩阵和关键变量散点关系，识别影响整机效率的主要运行参数。

第 6 章总结全文研究结论，归纳本文工作的主要认识，并对后续在线状态评估和模型扩展方向进行展望。
